\documentclass[10pt,a4paper]{amsart}
\usepackage[margin=19mm]{geometry}
\usepackage{amsmath,amssymb,mathtools}
\usepackage[scr=rsfs,cal=euler]{mathalfa}
\usepackage{xfrac}
\usepackage[unicode,hidelinks,bookmarks=false]{hyperref}
\newcommand{\ie}{i.e.\ }
\newcommand{\cf}{cf.\ }
\newcommand{\resp}{respectively\ }
\newcommand{\Subsection}{\subsection}
\providecommand{\nobreakdash}{\nobreak\hbox{-}}
\setlength{\emergencystretch}{2em}
\multlinegap0pt
\usepackage{bm}
\usepackage{bbm}
\usepackage{dsfont}
\usepackage{xspace}
\usepackage{cancel}
\usepackage{mathtools}
\usepackage{amsthm}
\makeatletter
\@ifundefined{proposition}{\newtheorem{proposition}{Proposition}}{}
\makeatother
\title[A unified parametric approach to the Erd\H{o}s--Straus conje]{A unified parametric approach to the Erd\H{o}s--Straus conjecture with explicit solutions for a set of integers of natural density one}
\author{Philemon Urbain Mballa}
\date{Source: arXiv:2602.20036v2 (CC BY 4.0)}
\begin{document}
\maketitle

\noindent\emph{Source and licence.} A unified parametric approach to the Erd\H{o}s--Straus conjecture with explicit solutions for a set of integers of natural density one. Version arXiv:2602.20036v2 (CC BY 4.0), distributed under the Creative Commons Attribution 4.0 International licence, as recorded in the arXiv licence metadata for this exact version and shown on the abstract page https://arxiv.org/abs/2602.20036v2. Full rights notice: licenses/arxiv-2602.20036-CC-BY-4.0.txt.\par\smallskip

\noindent\textit{Editorial scope.} Counted unit: the Proposition whose source label is \texttt{None} in the article (source0.tex lines 389--404, source label \texttt{None}), delivered together with the article's own complete original proof of it (source0.tex lines 405--478, one \texttt{proof} environment of 74 lines). No mathematics is written, repaired or re-derived: statement and proof are the article's own text line for line. Required context retained uncounted: none. Every definition, display and label of those lines is retained unchanged. Omitted from this package and not redistributed: 1--388, 479--813. Nothing inside a retained excerpt is omitted or altered: every retained body is delivered exactly as the article's lines are written, so no omission receipt of any kind accompanies this package. Citations: the retained excerpts carry no citation command at all, so no bibliography entry of the article is transcribed and no reference is redistributed. Reference adaptation: none was needed and none was made; every reference inside the delivered unit resolves inside it, so no unresolved reference or citation is produced. Printed slips kept literal: no printed word, symbol, hypothesis or number of the counted statement or of its proof was altered. Layout and class: the article's own class is \texttt{article} with packages including inputenc, fontenc, amsmath, amssymb, amsthm, hyperref; the wrapper is the lane's pinned \texttt{amsart} editorial wrapper at 10pt on A4, which restates the theorem-like environments the retained text needs and adds the article's own macro definitions that the retained text needs (none), with extra wrapper packages (bm, bbm, dsfont, xspace, cancel, mathtools). The delivered numbering is the unit's own. Assets: no retained span contains a figure environment, a graphics-inclusion command, a PDF-page inclusion, a table or a TikZ picture; no image file is copied into this package, the package declares no asset dependency and no image file is redistributed with it. Licence: version arXiv:2602.20036v2 of this article is distributed under the Creative Commons Attribution 4.0 International licence, as recorded in the arXiv licence metadata for this exact version and shown on the abstract page https://arxiv.org/abs/2602.20036v2. A full rights notice is delivered at licenses/arxiv-2602.20036-CC-BY-4.0.txt. Characters: every retained excerpt is pure ASCII, so the delivered engine, XeLaTeX with its default Latin Modern text font, reports no missing character.
  \begin{proposition}[Explicit symmetric solutions for a subfamily of \(n \equiv 1 \pmod{4}\)]
Let \(n \equiv 1 \pmod{4}\), written as \(n = 4r + 1\) with \(r \in \mathbb{N}^*\). 
Let \(b\) be an odd integer such that \(b \equiv 3 \pmod{4}\), and define
\[
x = \frac{n+b}{4}.
\]
If \(b \mid n\) (equivalently \(b \mid (4r+1)\)), then there exists an integer \(t \in \mathbb{N}^*\) such that
\(F_{x,t}^{(4)}(n) = 0\). This yields the symmetric solution
\[
y = z = t(4x-n) = t b,
\]
satisfying
\[
\frac{4}{n} = \frac{1}{x} + \frac{1}{y} + \frac{1}{z}.
\]
\end{proposition}

\begin{proof}
Since $n = 4r+1$ and $b \equiv 3 \pmod{4}$, we have
\[
n+b \equiv 1+3 \equiv 0 \pmod{4},
\]
hence
\[
x = \frac{n+b}{4} \in \mathbb{N}^*.
\]

We compute
\[
4x-n = 4\frac{(n+b)}{4} - n = n+b-n = b .
\]

From the definition

\(F_{x,t}^{(4)}(n)=t^2(4x-n)^2 - 2nxt\)

the condition \(F_{x,t}^{(4)}(n)=0\) gives
\[
t = \frac{2nx}{(4x-n)^2}, 
\]

Since $4x-n = b$, this becomes
\[
t = \frac{2nx}{b^2}.
\]

Substituting $n=4r+1$ and $x=\frac{(4r+1+b)}{4}$ yields
\[
t
= \frac{2(4r+1)\frac{4r+1+b}{4}}{b^2}
= \frac{(4r+1)(4r+1+b)}{2b^2}.
\]

Assume now that $b \mid (4r+1)$. Then there exists an integer $w \ge 1$ such that
\[
4r+1 = bw.
\]

Hence
\[
4r+1+b = bw + b = b(w+1),
\]
and therefore
\[
(4r+1)(4r+1+b) = b^2 w(w+1).
\]

Substituting into the expression of $t$, we obtain
\[
t = \frac{b^2 w(w+1)}{2b^2}
  = \frac{w(w+1)}{2}.
\]

Since $w$ and $w+1$ are consecutive integers, one of them is even.
Thus $w(w+1)$ is divisible by $2$, and consequently
\[
t \in \mathbb{N}^*.
\]

Therefore, for every $b \equiv 3 \pmod{4}$, the condition
\[
b \mid (4r+1)
\]
defines an infinite arithmetic progression in $r$, and hence in $n$.
This provides infinitely many explicit subfamilies in the class
$n \equiv 1 \pmod{4}$ yielding solutions
\[
y = z = t(4x-n)=tb
\]
that satisfy the Erd\H{o}s--Straus equation.
\end{proof}

\end{document}
